\chapter*{Abstract}
\addcontentsline{toc}{chapter}{Abstract}

Data-Science-Projekte sind durch hohe Unsicherheit, komplexe Datenabhängigkeiten 
und schwer interpretierbare Modelle gekennzeichnet, wodurch klassische Ansätze des 
Projekt- und Risikomanagements nur eingeschränkt geeignet sind. Insbesondere Risiken wie Datenbias, 
mangelnde Datenqualität, Modellinstabilität oder regulatorische Anforderungen erfordern 
eine spezifische und kontinuierliche Betrachtung entlang des gesamten Projektverlaufs.

Ziel dieser Arbeit ist die Entwicklung einer strukturierten Risikomanagement-Strategie 
für Data-Science-Projekte, die sich systematisch in das Data Science Process Model (DASC-PM) integrieren lässt. 
Hierzu werden zunächst grundlegende Begriffe, Risikotypen sowie relevante Standards und Methoden des 
Risikomanagements aufgearbeitet. Darauf aufbauend wird eine sogenannte „RM-Extension“ entwickelt, 
die das DASC-PM um phasenspezifische Risk Gates erweitert. Diese fungieren als Entscheidungs- und Kontrollpunkte, 
an denen Risiken mithilfe von Key Risk Indicators bewertet und gesteuert werden.

Die entwickelte Strategie wird im Rahmen von zwei Fallstudien exemplarisch angewendet, um 
typische Risiken in verschiedenen Projektphasen zu identifizieren, zu bewerten und geeignete 
Maßnahmen abzuleiten. Die Ergebnisse zeigen, dass eine systematische Integration von Risikomanagement 
in den Projektprozess dazu beitragen kann, Risiken frühzeitig zu erkennen, Transparenz zu erhöhen und 
fundierte Entscheidungen entlang des Projektverlaufs zu ermöglichen.

Die Arbeit verdeutlicht, dass ein phasenorientierter und indikatorbasierter Ansatz eine sinnvolle 
Ergänzung bestehender Prozessmodelle darstellt, weist jedoch auch auf Limitationen hinsichtlich 
der Generalisierbarkeit und der praktischen Umsetzung hin.
